\section{Introduction}
\label{sec:introduction}

Classical planning has become a cornerstone of automated decision-making in artificial intelligence, producing action sequences that transform an initial state into one that satisfies a goal specification. A planning problem takes as input a finite set of propositions, a set of operators with preconditions and effects over those propositions, an initial state, and a goal specification, and asks whether a finite sequence of applicable operators reaches a state in which the goal holds. Determining whether such a plan exists matters in practice. Planner runtime is wasted on unsolvable instances, model errors must be located and repaired before redeployment, and in safety-critical settings the absence of a plan signals that the operational specification needs revision. When a planning problem admits a solution, modern planners return a plan. If no solution exists, they typically report only that the problem is unsolvable. This binary failure report offers no account of why no plan exists, leaving users without the information needed to repair the problem specification or revise the modeling assumptions. Inconsistency measurement in propositional logic, surveyed by Thimm~\cite{Thimm2019} and by Grant and Hunter~\cite{Grant2023}, supplies a complementary line of research that quantifies the degree and structure of conflict in logical knowledge bases. Inconsistency measures assign numerical scores to knowledge bases and are characterized by rationality postulates that constrain how the score should respond to structural changes such as adding new formulas, enabling principled comparison and ranking of conflicts. Adaptations to related formalisms, including \ac{ASP} by Ulbricht et al.~\cite{Ulbricht2016} and temporal logic by Corea et al.~\cite{Corea2022}, show that the methodology transfers beyond its original setting, yet a systematic adaptation to classical planning remains open.

Existing diagnostic work on unsolvable planning problems addresses the gap only partially. Eriksson et al.~\cite{Eriksson2018} provide rule-based proof systems that certify unsolvability but return binary judgments, and Göbelbecker et al.~\cite{Goebelbecker2010} introduce ``excuses'' that prescribe minimal repairs without characterizing the underlying conflict structure. Neither approach delivers quantitative information about how unsolvability is structured within a problem. A natural shortcut is to encode a planning problem as a propositional knowledge base and apply a classical inconsistency measure directly, but this strategy fails. For instance, the contension measure $\mathcal{I}_c$ of Grant and Hunter~\cite{Grant2011}, which counts the minimum number of atoms that must take a conflicting truth value, returns a uniform value of one regardless of whether unsolvability arises from a missing operator, a pair of mutually exclusive goals, or an irreversible sequencing conflict, as shown in Section~\ref{sec:naive_encoding}. Planning conflicts manifest as reachability and achievability failures in a state-transition system rather than as direct logical contradictions, and any faithful measure must respect this planning-specific semantics.

The goal of this thesis is to adapt inconsistency measurement to classical planning in a way that operates directly on goals, operators, and reachable states rather than reducing to propositional encodings. Three planning-native diagnostic profiles comprising six measures in total are proposed. Each profile pairs a \emph{scope} measure counting affected goal propositions with a \emph{structure} measure counting structural artifacts such as unreachable propositions, mutex pairs, or sequencing-conflict pairs. The Unreachability Profile $\mathcal{I}_{\text{UR}}$ counts goal propositions absent from every reachable state together with the total number of unreachable propositions, revealing cascading dependency failures. The Goal Mutex Profile $\mathcal{I}_{\text{MX}}$ detects operator-induced mutual exclusions between goals that are individually achievable but never satisfiable in a common reachable state. The Goal Sequencing Conflict Profile $\mathcal{I}_{\text{GS}}$ identifies irreversible exclusions in which achieving one goal permanently blocks another. Rationality postulates specify properties, such as zero-iff-solvable and monotonicity under operator or goal additions, that a measure is expected to satisfy. Each profile is parameterized by a plan-length horizon $h$ that bounds the number of operator applications considered during reachability analysis, analyzed against planning-adapted versions of these classical postulates, and shown to be PSPACE-complete. The profiles are implemented in \ac{ASP}, using \textit{plasp}~3 by Dimopoulos et al.~\cite{Dimopoulos2019} for \ac{PDDL} translation and Clingo by Gebser et al.~\cite{Gebser2012} for state-space exploration under brave reasoning. Brave reasoning is a query mode in which an atom holds if it appears in at least one answer set, so each answer set witnesses one execution trace. The implementation is evaluated on the \ac{IPC}~2016 Unsolvability Competition benchmarks by Muise and Lipovetzky~\cite{Muise2016} and the unsolvable benchmark collection by Eriksson~\cite{Eriksson2019}.

\subsection{Problem Statement}
\label{sec:problem_statement}

Current planning systems provide minimal diagnostic information when they encounter unsolvable problems. Users lack insight into why a problem has no plan and have no systematic means to distinguish or quantify structurally different cau\-ses of unsolvability. Consequently, identifying minimal changes that would restore solvability reduces to trial-and-error over the problem specification, since different types of unsolvable planning problems are neither named nor measured. This thesis supplies the missing quantitative diagnostic layer through three planning-native diagnostic profiles comprising six measures in total, targeting unreachability, operator-induced mutex, and irreversible sequencing conflicts.

\subsection{Research Questions}
\label{sec:research_questions}

The central research question of this thesis is:

\begin{quote}
    How can inconsistency measures from propositional logic be meaningfully adapted to classical planning to provide diagnostic value for unsolvable planning problems?
\end{quote}

This central question is decomposed into three research questions, each addressed in a dedicated section of this thesis.

\paragraph{Research Question~1: Theoretical Adaptation.}
\begin{quote}
    Which inconsistency measures from propositional logic can be adapted to classical planning, and what postulates do they retain?
\end{quote}
A direct application of the contension measure to a propositional encoding of a planning problem is shown to be uninformative, motivating the development of planning-native measures. This question is answered in Section~\ref{sec:theoretical_adaptation}.

\paragraph{Research Question~2: Implementation.}
\begin{quote}
    How can the adapted measures be efficiently computed for planning problems using Answer Set Programming?
\end{quote}
This question is answered in Section~\ref{sec:implementation}, which describes the three-phase pipeline and its automated test suite.

\paragraph{Research Question~3: Practical Evaluation.}
\begin{quote}
    Do the adapted measures provide practical diagnostic value for understanding and analyzing unsolvable planning problems?
\end{quote}
This question is answered in Section~\ref{sec:evaluation}, where seven compatible domains are analyzed across 158 instances.

\subsection{Contributions}
\label{sec:contributions}

The contributions of this thesis are organized along the three research questions.

\paragraph{Theoretical Adaptation.}
The thesis demonstrates that naive propositional encodings of planning problems yield uninformative inconsistency measure values, with the contension measure returning one regardless of structural cause. Three plan\-ning-native diagnostic profiles comprising six measures in total are proposed and formally defined as horizon-parameterized profiles consisting of a scope measure and a structure measure: the Unreachability Profile $\mathcal{I}_{\text{UR}}$, the Goal Mutex Profile $\mathcal{I}_{\text{MX}}$, and the Goal Sequencing Conflict Profile $\mathcal{I}_{\text{GS}}$. This thesis develops planning-adapted versions of the classical rationality postulates, namely Planning Consistency, Operator Monotonicity, and Goal Monotonicity, and analyzes all six measures against them. The containment relation between the measures is established, with sequencing conflicts implying mutex and unreachable goals excluded from both conflict analyses. The decision problems $\text{UPPER}$, $\text{LOWER}$, and $\text{EXACT}$ for all three measures are shown to be PSPACE-complete, yielding a uniform complexity landscape that contrasts with the spread of complexity classes across measures previously observed for propositional and temporal logic adaptations.

\paragraph{Implementation.}
The measures are realized in an end-to-end \ac{ASP} pipeline that accepts standard \ac{PDDL} files and returns the full diagnostic profile at a user-specified horizon. The central methodological contribution is the use of brave reasoning as a decision procedure for mutex and sequencing conflicts. Different answer sets correspond to different execution traces, and the absence of a witness atom under brave reasoning constitutes a universal proof that the corresponding condition is unsatisfiable across all traces. Correctness is verified through 80 automated tests covering scenario-level validation, hierarchy relationships, library-\ac{API} contracts, and end-to-end \ac{PDDL} processing, with all eleven test scenarios producing exact matches with their theoretically derived profiles.

\paragraph{Practical Evaluation.}
An empirical study on the \ac{IPC}~2016 Unsolvability Competition and the Eriksson benchmark collection establishes the practical diagnostic value of the measures. Of 158 unsolvable instances across seven compatible domains, 68 (43\%) complete within the 120-second time limit at horizon $h = 20$, with the resulting profiles distinguishing unreachability-dominated from mutex-dominated and sequencing-dominated domains. The measures provide three levels of diagnostic granularity: category classification, scope at the goal level, and structure at the relationship level. A horizon sensitivity analysis exposes a soundness-coverage tradeoff and motivates the default horizon of $h = 20$, which avoids observed false positives. The \textit{3unsat} domain supplies the first empirical evidence of a pairwise-analysis blind spot, where instances with combinatorial unsatisfiability return zero profiles despite being unsolvable.

\subsection{Thesis Structure}
\label{sec:thesis_structure}

Section~\ref{sec:foundations} establishes the theoretical foundations, covering inconsistency measurement for propositional knowledge bases, classical planning formalisms, computational complexity, and \ac{ASP}. The section concludes by identifying the challenges that distinguish planning from static knowledge bases and motivate a planning-native adaptation.

Section~\ref{sec:related_work} surveys the inconsistency measurement literature and positions the thesis against existing work on planning diagnostics, including dead-state certificates \cite{Eriksson2018}, repair-based excuses \cite{Goebelbecker2010}, and explanation methods for unsolvable tasks \cite{Sreedharan2019}. It also discusses adaptations of inconsistency measurement to non-propositional domains, which serve as methodological precedents.

Section~\ref{sec:theoretical_adaptation} develops the three planning-native measures. It opens with a structural taxonomy of unsolvable planning problems, demonstrates the limitations of naive propositional encodings on six running test scenarios, introduces horizon-bounded reachability, and formally defines the Unreachability, Goal Mutex, and Goal Sequencing Conflict Profiles. The section then analyzes the measures against rationality postulates, establishes formal relationships between them, and proves uniform PSPACE-completeness of their decision problems.

Section~\ref{sec:implementation} presents the \ac{ASP}-based implementation. It describes the three-phase pipeline, the adoption of \textit{plasp} for \ac{PDDL} translation, the reachability and witness-collection encodings used under Clingo brave reasoning, the Python module that extracts measure values, and the automated test suite that validates the pipeline.

Section~\ref{sec:evaluation} reports the experimental evaluation. It describes the methodology, pre\-sents per-domain results across the seven compatible domains, discusses case studies that illustrate the diagnostic capabilities and the pairwise-analysis limitation, analyzes horizon sensitivity, and examines computational performance.

Section~\ref{sec:conclusion} summarizes the contributions along the three research questions, discusses limitations, provides a closing assessment of the methodological insight that planning conflicts manifest as reachability and achievability failures rather than direct logical contradictions, and identifies directions for future work.
